% =====================================================================
\section{Conclusion}\label{sec:conclusion}
% =====================================================================
Physical faithfulness in generated video requires three things that are usually
studied apart: a way to \emph{steer} what gets generated, a way to \emph{measure}
whether physical obligations are met, and \emph{ground truth} against which
measurements can be scored. This report contributes one artifact toward each. For
steering, we build a simulation-driven controllable generator
(MuJoCo\,$\rightarrow$\,Wan~2.2-VACE), a first step toward the control channel
through which a critic might one day feed corrections back into generation. For
ground truth, we assemble a benchmark of human-annotated, localized defects, so
that a critic can be scored on whether it finds the specific fault rather than on
aggregate preference. And for measurement, we present \PhyReAct{}, an auditable
multi-tool critic that grounds the reason--act--observe loop in physical
verification: reasoning compiles the prompt into typed physical checks, frozen
expert operators are called only to return measurements over the scope those
checks declare, and a deterministic roll-up folds the observations into a
three-valued verdict (\supported{}/\contradicted{}/\unknownv{}, surfaced as
\emph{plausible}, \emph{implausible}, or \abstain{}) with a human-traceable
evidence chain.

On a development core of $149$ clips carrying $304$ human-marked defects,
\PhyReAct{} accounts for $228$, against $164$ for a published question-decomposition
evaluator reading the same clips and checking the same claims with the same served
model. Prompting the same backbone monolithically reaches $222$, so the count alone does not
separate the two. Every one of \PhyReAct{}'s $228$ arrives with the number that forced it, and
a reader can audit any single verdict instead of accepting an aggregate.

That auditability is also what makes the evidence usable downstream. A contradicted
claim names what failed, when it failed, and on what measurement, which is precisely
the signal a generator needs to regenerate a clip locally rather than resample it
whole. It is equally the substrate for an operator library that grows as new
measurement procedures are distilled, and for improvement that edits a readable state
rather than a weight. Closing the critic-to-generator loop over that signal is the
next milestone. \Cref{sec:framework} casts these components as a single state-optimization
framework, and \cref{sec:roadmap} lays out the phased plan that carries it there.

% =====================================================================
